\documentclass[10pt]{article}
\usepackage[margin=2.0cm]{geometry}
\usepackage{multicol}
\usepackage[compact]{titlesec}
\usepackage{enumitem}
\usepackage[hidelinks]{hyperref}
\setlength{\parskip}{0.5pt}
\setlength{\multicolsep}{3pt}
\titlespacing{\section}{0pt}{*0.8}{*0.3}
\setlist[itemize]{nosep, leftmargin=1.3em}

\title{\textbf{Measuring and Modeling Stage Placement in Hybrid Retrieval-Augmented Generation across Edge, LAN, and WAN}}
\author{M. Abdullah Ali \and M. Abdullah Aamir\\
\small Department of Data Science, FAST-NUCES, Islamabad \quad---\quad Proposed supervisor: Dr.~Adil Majeed}
\date{September 14, 2026}

\begin{document}
\maketitle
\vspace{-9mm}

\section{Background and Current State}
We have built and deployed a working \emph{hybrid} Retrieval-Augmented Generation (RAG) system:
sparse keyword retrieval (BM25 via Tantivy), dense semantic retrieval (BGE-M3 embeddings searched
in Qdrant), reciprocal-rank fusion, and LLM generation (Llama-3-8B-Instruct, AWQ 4-bit, served by
vLLM~\cite{vllm,bgem3}). The corpus is the full MS MARCO passage collection~\cite{msmarco}:
\textbf{8{,}841{,}823 passages} embedded at 1024 dimensions, HNSW-indexed in Qdrant, ingested at
roughly 538 passages per second on an RTX 3090. The system is deployed on two LAN machines
connected over gRPC (WireGuard): a gateway node (embedding, dense+sparse retrieval, fusion) on an
Intel i7-10750H
laptop with GTX 1660~Ti, and a generation node (document lookup, vLLM) on an RTX 3090 host. An
earlier three-node version placed a third, low-power laptop (Intel CPU, 16~GB RAM) at the network
edge as entry point and sparse-retrieval engine, across a city-wide link treated as the WAN leg.

\section{A Failed Hypothesis That Points at a Real Gap}
The original design hypothesized a 30--40\% TTFT (time-to-first-token) reduction by running
edge-side sparse retrieval in parallel with transit of the query toward the other nodes:
$T_{\mathrm{retrieval}} = \max(T_{\mathrm{sparse}},\; T_{\mathrm{WAN}} + T_{\mathrm{dense}})$.
Our instrumented three-node runs showed nothing like it. The surviving 100-query log records
sparse retrieval at 0.6--10~ms against TTFT of 51--227~ms: the stage we chose to distribute was
one to two orders of magnitude cheaper than the pipeline it was meant to accelerate. Since
parallel overlap can save at most
$\min(T_{\mathrm{sparse}},\; T_{\mathrm{WAN}} + T_{\mathrm{dense}})$, no link latency could have
delivered the target --- a short WAN leg offered little to hide, and a long one would only have
slowed the dense branch. One caveat: we can no longer reconstruct where the machines physically
were during those logged runs, so we treat the 5--15~ms round-trip from our mid-evaluation report
as descriptive, not tied to this log; the new campaign will record placement and network metadata
with every measurement. The three-node setup
itself ended for operational reasons, not hardware: the edge laptop belonged to a groupmate who was
often unavailable, so we consolidated its duties into the gateway and the two-node system worked
fine.

That failure motivates this proposal: placement by intuition silently costs performance, the
stages worth distributing are the \emph{expensive} ones (dense retrieval, document handling,
generation --- not sparse lookup), and topology change is a normal event in real deployments.
Splitting is often not a choice: large models need GPUs most machines do not have, data-residency
rules can force retrieval and generation onto different hosts, and cost pushes work toward
whatever hardware exists. Yet the literature, as far as we can find, cannot
answer: \emph{where should each stage of a hybrid RAG pipeline run, under which network
conditions, and what exactly do we lose when the topology degrades?}

\section{The Gap in the Literature}
DNN partitioning between device and cloud is well studied at the \emph{layer} level
(Neurosurgeon~\cite{neurosurgeon}, PipeEdge~\cite{pipeedge}); LLM serving splits model \emph{phases}
across machines (Splitwise~\cite{splitwise}). For RAG specifically, existing work optimizes KV
caching of retrieved knowledge (RAGCache~\cite{ragcache}, TurboRAG~\cite{turborag}),
serving-system optimization (RAGO~\cite{rago}, METIS~\cite{metis}), index management on a single
host (EdgeRAG~\cite{edgerag}, VectorLiteRAG~\cite{vectorliterag}), or data distribution
across edge/cloud tiers (EACO-RAG~\cite{eacorag}, DGRAG~\cite{dgrag}). We could not find a
measurement study and placement policy for \emph{which pipeline stage runs on which machine} in a
multi-node hybrid RAG deployment: the five stages (sparse, dense, fusion, hydration, generation)
can each sit on a different node, and the right answer changes with RTT, bandwidth, and node
capability. Node loss (what happens to latency and answer quality when a node drops out and roles
consolidate) appears entirely unmeasured for RAG.

\section{Research Questions}
\begin{itemize}
  \item \textbf{RQ1 (Placement).} Which placement for each stage minimizes TTFT and end-to-end
  latency under a given network regime (RTT, bandwidth), and where does the best placement
  change?
  \item \textbf{RQ2 (Degradation).} When a node drops out --- hardware failure or simple
  unavailability (our real three-to-two consolidation) --- what is the cost in latency and
  retrieval quality, and how fast can the system adapt?
  \item \textbf{RQ3 (Policy).} Can a simple analytical cost model --- per-stage compute times,
  payload sizes, network terms --- predict the best placement within a stated tolerance (e.g.\
  TTFT within 15\% of measured), and can the running system use it to pick placement per network
  regime without the full experiment matrix?
\end{itemize}

\section{Method}
\textbf{Testbed.} The existing heterogeneous nodes (RTX 3090 host, GTX 1660~Ti laptop, revived
low-power laptop) over WireGuard; the full 8.84M-passage index is already in place.
\textbf{Instrumentation.} We will extend the existing per-stage timing logs into a structured
benchmark harness (per-stage milliseconds, TTFT, decode throughput, true stream end-to-end) and add
retrieval-quality evaluation against MS MARCO qrels (recall@$k$) plus an answer-level check
(exact match), so topology-induced damage shows up in generated answers and not only in retrieval;
the sparse index will be scaled from a small test index to the full corpus.
\textbf{Network control.} Today the only delay knob is an application-level header; we will add
kernel-level per-link shaping (Linux netem on each WireGuard interface: delay and bandwidth,
0/10/40/80~ms RTT) so regimes shape every hop, and re-validate key configurations over the genuine
WAN link. \textbf{Placement matrix.} Stage-to-node assignment becomes a configuration switch
(fusion on gateway vs.\ generator; hydration local vs.\ remote; sparse retrieval at the edge
vs.\ LAN), reconstructing the archived three-node code paths. \textbf{Campaign.} A factorial sweep
over placement $\times$ network regime $\times$ retrieval depth with repeats and percentile
statistics; a degradation experiment kills nodes mid-load to measure failover cost; the RQ3 cost
model is fit on half the configurations, validated on the held-out half, and then run live, letting
the model pick placement per regime against the fixed ones. \textbf{Limitations.} Three machines,
one per capability tier: results characterize this class of testbed, not a fleet; netem regimes
bound generality.

\section{Expected Contributions}
(1) To our knowledge, the first factorial measurement of stage-to-node placement for hybrid RAG on
commodity
heterogeneous hardware under controlled WAN regimes; (2) a measurement of what topology
degradation actually costs, grounded in our own three-to-two-node consolidation; (3) a cost-model
placement policy, validated on held-out configurations and exercised live; (4) an open,
reproducible testbed with all artifacts.

\section{Timeline and Target Venues}
Weeks 1--2: harness, full-corpus sparse index, quality evaluation. Week 3: placement matrix and
network emulation. Weeks 4--5: campaigns, cost model, analysis. Weeks 6--8: writing, internal
review, submission. If time runs short we cut in order: live policy validation, extra network
regimes --- never the degradation study or the quality metrics. We target Q1 journals with rolling
submission: \emph{Future Generation Computer
Systems}, \emph{Computer Networks}, or \emph{IEEE Transactions on Cloud Computing} (all currently
Q1 in SCImago/JCR rankings). A conference extension (e.g.\ IEEE ICDCS 2027, call expected around
December 2026) remains possible for the policy contribution.

\begin{multicols}{2}
\begin{thebibliography}{9}\scriptsize
\bibitem{neurosurgeon} Y.~Kang et al., ``Neurosurgeon: Collaborative Intelligence Between the Cloud
and Mobile Edge,'' \emph{Proc.\ ASPLOS}, 2017.
\bibitem{pipeedge} Y.~Hu et al., ``PipeEdge: Pipeline Parallelism for Large-Scale Model Inference on
Heterogeneous Edge Devices,'' \emph{Proc.\ DSD}, 2022.
\bibitem{splitwise} P.~Patel et al., ``Splitwise: Efficient Generative LLM Inference Using Phase
Splitting,'' \emph{Proc.\ ISCA}, 2024.
\bibitem{ragcache} C.~Jin et al., ``RAGCache: Efficient Knowledge Caching for Retrieval-Augmented
Generation,'' \emph{ACM Trans.\ Comput.\ Syst.}, 2026.
\bibitem{turborag} S.~Lu et al., ``TurboRAG: Accelerating Retrieval-Augmented Generation with
Precomputed KV Caches for Chunked Text,'' \emph{Proc.\ EMNLP}, 2025.
\bibitem{rago} W.~Jiang et al., ``RAGO: Systematic Performance Optimization for Retrieval-Augmented
Generation Serving,'' \emph{Proc.\ ASPLOS}, 2025.
\bibitem{metis} S.~Ray et al., ``METIS: Fast Quality-Aware RAG Systems with Configuration
Adaptation,'' \emph{Proc.\ SOSP}, 2025.
\bibitem{edgerag} K.~Seemakhupt, S.~Liu, and S.~Khan, ``EdgeRAG: Online-Indexed RAG for Edge
Devices,'' arXiv:2412.21023, 2024.
\bibitem{vectorliterag} J.~Kim and D.~Mahajan, ``VectorLiteRAG: Latency-Aware and Fine-Grained
Resource Partitioning for Efficient RAG,'' arXiv:2504.08930, 2025.
\bibitem{eacorag} J.~Li et al., ``EACO-RAG: Towards Distributed Tiered LLM Deployment Using
Edge-Assisted and Collaborative RAG with Adaptive Knowledge Update,'' arXiv:2410.20299, 2024.
\bibitem{dgrag} W.~Zhou, Y.~Yan, and Q.~Yang, ``DGRAG: Distributed Graph-based Retrieval-Augmented
Generation in Edge-Cloud Systems,'' arXiv:2505.19847, 2025.
\bibitem{msmarco} P.~Bajaj et al., ``MS MARCO: A Human Generated MAchine Reading COmprehension
Dataset,'' arXiv:1611.09268, 2016.
\bibitem{bgem3} J.~Chen et al., ``M3-Embedding: Multi-Linguality, Multi-Functionality,
Multi-Granularity Text Embeddings Through Self-Knowledge Distillation,'' arXiv:2402.03216, 2024.
\bibitem{vllm} W.~Kwon et al., ``Efficient Memory Management for Large Language Model Serving with
PagedAttention,'' \emph{Proc.\ SOSP}, 2023.
\end{thebibliography}
\end{multicols}

\end{document}
